\section{Direct confinement}
\label{sec:direct_confinement}

The trace identity turns information about \(\AX\) into information
about where the mass of \(\mu\) lies.  The qualitative statement uses
\(\AX=0\); the quantitative one uses only the size of \(\trace\AX\).

\begin{theorem}[Separation--confinement]
\label{thm:duality_confinement:separation-confinement}
Let \(\IOL\) be an involutive object ledger with
\(\int_X\|\psim\|^2\,d\mu<\infty\), and suppose \(\AX=0\).  Then
\(\psim(x)=0\) for \(\mu\)-almost every \(x\).  If moreover \(\psim\)
is qualitatively separating, then \(Jx=x\) for \(\mu\)-almost every
\(x\); that is, \(X\setminus\Fix(J)\) is \(\mu\)-null.  In the finite
weighted case, every object of positive weight is fixed by \(J\).
\end{theorem}

\begin{proof}
By \Cref{lem:duality_confinement:trace-identity},
\(\int_X\|\psim\|^2\,d\mu=\trace\AX=0\).  The integrand is nonnegative
and measurable, so \(\|\psim(x)\|=0\) for \(\mu\)-almost every \(x\).
If \(\psim\) is qualitatively separating, the set of \(x\) at which the
separation equivalence fails is also \(\mu\)-null; outside the union of
the two null sets, \(\psim(x)=0\) and hence \(Jx=x\).  In the finite
case the null sets contain no object of positive weight.
\end{proof}

\begin{theorem}[Quantitative confinement]
\label{thm:duality_confinement:quantitative-confinement}
Let \(\IOL\) be an involutive object ledger on a metric space \(X\),
with \(\int_X\|\psim\|^2\,d\mu<\infty\) and with \(\psim\)
quantitatively separating with modulus \(m\).  Then, for every
\(\varepsilon>0\),
\[
  \mu^*\bigl(\{x\in X:\dist(x,\Fix(J))\geq\varepsilon\}\bigr)
  \leq
  \frac{\trace\AX}{m(\varepsilon)^2},
\]
where \(\mu^*\) is the outer measure of \(\mu\); it equals \(\mu\) on
measurable sets.  The same bound holds with \(\dist(\cdot,\Fix(J))\)
replaced by any function \(\delta:X\to[0,\infty]\) such that
\(\delta(x)\geq\varepsilon\) implies \(\|\psim(x)\|\geq m(\varepsilon)\).
\end{theorem}

\begin{proof}
Fix \(\varepsilon>0\) and put
\(E_\varepsilon=\{x:\dist(x,\Fix(J))\geq\varepsilon\}\) and
\(F_\varepsilon=\{x:\|\psim(x)\|\geq m(\varepsilon)\}\).  Quantitative
separation gives \(E_\varepsilon\subseteq F_\varepsilon\), and
\(F_\varepsilon\) is measurable because \(\psim\) is strongly
measurable (up to a \(\mu\)-null set if \(\psi\) is only almost
everywhere strongly measurable, which does not affect the bound).  Markov's inequality \citep{Folland1999} and
the trace identity give
\[
  m(\varepsilon)^2\,\mu(F_\varepsilon)
  \leq\int_{F_\varepsilon}\|\psim\|^2\,d\mu
  \leq\int_X\|\psim\|^2\,d\mu=\trace\AX .
\]
Since \(m(\varepsilon)>0\), dividing and using
\(\mu^*(E_\varepsilon)\leq\mu(F_\varepsilon)\) gives the bound.  The
proof used the metric only through the inclusion
\(E_\varepsilon\subseteq F_\varepsilon\), which holds for any
\(\delta\) as in the statement.
\end{proof}

\begin{example}[Complex conjugation, continued]
In the example of \Cref{sec:involutive_ledger},
\(\Pminus(u,v)=(0,v)\), so \(\psim(z)=(0,\operatorname{Im}z)\) and
\[
  \AX=\Bigl(\int_X(\operatorname{Im}z)^2\,d\mu(z)\Bigr)
  \begin{pmatrix}0&0\\0&1\end{pmatrix},
\]
provided the integral is finite.  If \(X\) is a disc centred on the
real axis, then \(\operatorname{Re}z\in X\) for every \(z\in X\), so
\(\dist(z,X\cap\mathbb R)=|\operatorname{Im}z|\); the readout is then
quantitatively separating with \(m(\varepsilon)=\varepsilon\), and
\Cref{thm:duality_confinement:quantitative-confinement} becomes the
familiar Markov--Chebyshev bound
\(\mu\{|\operatorname{Im}z|\geq\varepsilon\}\leq
\varepsilon^{-2}\int_X(\operatorname{Im}z)^2\,d\mu\).
\end{example}

The confinement theorem of \Cref{sec:master_theorem} supplies the
missing input \(\AX=0\): it derives it from domination by operators of
vanishing trace, after which
\Cref{thm:duality_confinement:separation-confinement} applies.  The
quantitative theorem is useful when only an upper bound on
\(\trace\AX\) is available, as in
\Cref{prop:duality_confinement:optimized-trace-budget}.
